\subsection{K4/FOUR analytic spine}

\begin{theorem}[The K4/FOUR chromatic obstruction carries the finite analytic spine]\label{thm:k4-four-analytic-spine}\lean{Obsmetrologia.Flow.k4_four_analytic_spine}\leanok\uses{thm:k4-group, thm:k4-proper-colouring-count, thm:realized-tile-adjacency-chromatic-numbers, thm:tile-rayleigh-bottom, thm:two-site-gap, thm:three-site-gap-lower, thm:tile4-irrational-gap, thm:tile7-gap-bracket, thm:finite-predictor-phase-type, thm:finite-diffraction-mass}
The complete FOUR obstruction is not only a colouring lower bound.  The same
finite carrier now has a bundled analytic readout: the \(K_4\) action has four
involutive symmetries acting by composition on the charges; king adjacency is
FOUR-colourable and not three-colourable; \(K_4\) has \(24\) proper
FOUR-colourings; the chromatic tile has Rayleigh bottom \(0\), zero energy
exactly on the harmonic kernel, and the \(H^0/H^1\) index balance; the two- and
three-site gaps have exact finite witnesses; the four-site quadratic gap and
seven-site cubic bracket are present; and the finite Fourier-side mass and
phase readings agree with the predictor-level Pisot/Salem split.

The theorem is a finite spine theorem.  It does not construct the general
\(N\)-site spectral formula, the \(N=7\) Bishop real root, infinite-volume
diffraction, or an analytic Lebesgue decomposition.
\end{theorem}
